\chapter{Modèle de données}\label{chap:donnees}

\section{Schéma relationnel}
Le schéma (\code{schema.sql}) est compatible H2 et PostgreSQL : types standard, clés générées par
\code{GENERATED BY DEFAULT AS IDENTITY}. Les références entre domaines (par exemple \code{candidature.offre\_id})
sont volontairement des \textbf{identifiants logiques sans clé étrangère} : chaque service étant propriétaire de ses
données, l'intégrité référentielle inter-domaines est garantie par les vérifications du processus, comme elle le
serait avec des bases physiquement séparées.

\begin{figure}[htbp]
\centering
\begin{tikzpicture}[font=\sffamily\scriptsize, node distance=6mm and 9mm,
  ent/.style={draw=black!45, rounded corners=2pt, align=left, fill=#1, inner sep=4pt, text width=3.5cm},
  lien/.style={black!55, thick}, logique/.style={black!45, dashed, thick}]
  \node[ent=soap] (etu) {\textbf{etudiant}\\[2pt]\underline{matricule}\\nom, prenom, email\\telephone, adresse\\filiere, niveau, competences\\statut, date\_inscription};
  \node[ent=soap, right=of etu] (ent) {\textbf{entreprise}\\[2pt]\underline{id}\\raison\_sociale, secteur\_activite\\adresse, ville, email, telephone\\contact\_nom, contact\_email…\\referencee, date\_referencement};
  \node[ent=soap, right=of ent] (off) {\textbf{offre}\\[2pt]\underline{id}, entreprise\_id\\titre, description\\competences\_requises\\duree\_mois, date\_limite\\nombre\_places, places\_disponibles\\filieres\_ciblees, type\_stage\\statut, date\_publication};
  \node[ent=rest, below=14mm of etu] (cand) {\textbf{candidature}\\[2pt]\underline{id}, matricule, offre\_id\\entreprise\_id, nom\_etudiant\\titre\_offre, lettre\_motivation\\cv\_url, score\_compatibilite\\statut, motif\\process\_instance\_id\\date\_depot, date\_mise\_a\_jour};
  \node[ent=rest, below=of cand] (hist) {\textbf{candidature\_historique}\\[2pt]\underline{id}, candidature\_id\\ancien\_statut, nouveau\_statut\\commentaire, auteur, date};
  \node[ent=rest, right=of cand] (notif) {\textbf{notification}\\[2pt]\underline{id}, type\\destinataire\_type, destinataire\_id\\email, sujet, message\\canal, statut, lue\\candidature\_id, offre\_id\\date\_envoi};
  \node[ent=rest, right=of notif] (stage) {\textbf{stage}\\[2pt]\underline{id}, candidature\_id (unique)\\matricule, offre\_id, entreprise\_id\\sujet, date\_debut, date\_fin\\tuteur\_entreprise\\encadrant\_academique, statut};
  \node[ent=rest, below=of stage, xshift=-22mm] (cr) {\textbf{compte\_rendu}\\[2pt]\underline{id}, stage\_id\\periode, contenu\\auteur, date\_depot};
  \node[ent=rest, right=4mm of cr] (eval) {\textbf{evaluation}\\[2pt]\underline{id}, stage\_id, type\\evaluateur, note /20\\appreciation, date};
  \draw[logique] (ent) -- node[above]{1..n} (off);
  \draw[logique] (etu) -- node[left]{1..n} (cand);
  \draw[logique] (off.south) -- ++(0,-6mm) -| node[pos=0.25, above]{1..n} ([xshift=8mm]cand.north);
  \draw[lien] (cand) -- node[left]{1..n} (hist);
  \draw[logique] (cand) -- node[above]{0..1} (stage.west |- cand.east);
  \draw[lien] (stage.south) -- ++(0,-3mm) -| (cr.north);
  \draw[lien] (stage.south) -- ++(0,-3mm) -| (eval.north);
\end{tikzpicture}
\caption{Modèle de données (bleu : services SOAP, vert : services REST ; pointillés : références logiques entre domaines)}
\end{figure}

\section{Énumérations}
\begin{table}[htbp]
\centering
\caption{Valeurs énumérées}
\begin{tabularx}{\textwidth}{@{}lY@{}}
\toprule
\textbf{Attribut} & \textbf{Valeurs} \\
\midrule
Statut étudiant & \code{ACTIF}, \code{SUSPENDU}, \code{EN\_STAGE}, \code{DIPLOME} \\
Année du cycle ingénieur & \code{INE1}, \code{INE2}, \code{INE3} \\
Type de stage (niveaux ouverts) & \code{OBSERVATION} (INE1), \code{OUVRIER\_PFA} (INE1, INE2), \code{PFE} (INE3) \\
Statut d'offre & \code{ACTIVE}, \code{SUSPENDUE}, \code{POURVUE}, \code{CLOTUREE} \\
Statut de candidature & \code{DEPOSEE}, \code{EN\_COURS\_EVALUATION}, \code{ACCEPTEE}, \code{REFUSEE}, \code{ANNULEE} \\
Statut de stage & \code{PLANIFIE}, \code{EN\_COURS}, \code{TERMINE}, \code{INTERROMPU} \\
Type d'évaluation & \code{MI\_PARCOURS}, \code{FINALE} \\
Type de notification & \code{CONFIRMATION\_DEPOT}, \code{RECEPTION\_CANDIDATURE}, \code{REJET\_CANDIDATURE}, \code{ACCEPTATION},
  \code{REFUS}, \code{DEMANDE\_COMPLEMENT}, \code{RAPPEL\_DATE\_LIMITE}, \code{RELANCE\_DECISION}, \code{INCIDENT\_TRAITEMENT},
  \code{ANNULATION}, \code{LIBRE} \\
Rôle utilisateur & \code{RESPONSABLE}, \code{ETUDIANT}, \code{ENTREPRISE}, \code{SYSTEME} \\
\bottomrule
\end{tabularx}
\end{table}

\section{Jeu de données de démonstration}
Le jeu de démonstration reprend le contexte de l'INPT : 8 élèves ingénieurs des filières
\textbf{AMOA, ASEDS, DATA, ICCN, SESNum, SUD} et \textbf{Cloud \& IoT} (niveaux notés comme à l'INPT :
INE1, INE2 et INE3 pour les trois années du cycle ingénieur), dont une élève
suspendue ; 6 entreprises fictives implantées à Casablanca, Rabat, Agadir et Marrakech (dont une non référencée) ;
12 offres couvrant les trois types de stage (observation, ouvrier / PFA, PFE), dont une à date limite dépassée et une
d'entreprise non référencée ; chaque année d'études dispose ainsi d'offres qui lui sont ouvertes. Chaque compte reçoit une clé d'API de
démonstration (annexe~\ref{annexe:comptes}). Ce jeu couvre volontairement tous les cas de rejet du processus.
